\chapter{RESEARCH METHODOLOGY}
\label{ch:methodology}

\section{Introduction}
\label{sec:meth_intro}

This chapter sets out the methodology by which the research questions
articulated in Chapter~\ref{ch:introduction} are addressed and the
hypotheses are tested.  It elaborates the research philosophy and
design (Section~\ref{sec:design}); the populations and sampling
strategies for the two primary instruments
(Section~\ref{sec:sampling}); the construction, scaling and
psychometric properties of the instruments
(Sections~\ref{sec:instruments}--\ref{sec:reliability}); the
hypothesis model (Section~\ref{sec:hyp_model}); the pilot study
(Section~\ref{sec:pilot}); the data-collection procedure
(Section~\ref{sec:datacoll}); the data-analysis plan
(Section~\ref{sec:analysis}); and the ethical safeguards adopted
(Section~\ref{sec:ethics}).

\section{Research Philosophy and Paradigm}
\label{sec:philosophy}

The study adopts a \emph{pragmatist} research philosophy
\citep{creswell2018,morgan2014}.  Pragmatism rejects the false binary
of positivism and constructivism and instead privileges the research
question as the primary determinant of method choice
\citep{tashakkori2010,johnson2007}.  For the present study, the
research questions span objective behavioural change (RQ1, RQ2),
structural-economic reconfiguration (RQ3) and interpretive
sense-making by industry actors (RQ4); a single epistemological
commitment is therefore inadequate.  Quantitative methods are
appropriate for measuring the magnitude and statistical
significance of behavioural change in a 400-respondent viewer
sample; qualitative methods are appropriate for capturing the
strategic intent, narrative framings and informal market-conventions
of the 50 industry specialists.  The pragmatist orientation legitimises
the integration.

\section{Research Design}
\label{sec:design}

The thesis adopts a \textbf{sequential explanatory mixed-methods design}
\citep{creswell2018,tashakkori2010}, organised in three phases
(Figure~\ref{fig:hypothesis} presents the hypothesis model that
emerges from the design):

\begin{enumerate}[itemsep=2pt]
\item \textbf{Phase~1 -- Quantitative}.  A structured questionnaire
      administered to 400 Tamil-film viewers (October 2024 -- February
      2025).  Phase~1 tests hypotheses H1--H5
      (Section~\ref{sec:hyp_model}).
\item \textbf{Phase~2 -- Mixed quantitative + qualitative}.  A
      structured-cum-open instrument administered to 50 industry
      specialists, supplemented by 28 in-depth interviews of
      45--90 minutes (March 2025 -- May 2025).  Phase~2 tests
      hypotheses H6--H8 and provides explanatory depth for Phase~1.
\item \textbf{Phase~3 -- Triangulation and synthesis}.  Phase-1 and
      Phase-2 findings are cross-validated against secondary trade
      data (FICCI--EY, Ormax, KPMG, BCG, Statista, MPA, trade press)
      and synthesised into the Hybrid-Equilibrium framework
      (Chapter~\ref{ch:conclusion}).
\end{enumerate}

This design is well-suited to research where one phase informs the
interpretation of the other and where neither method alone could
answer the research questions.  The use of secondary data as a
third leg of triangulation strengthens external validity.

\section{Population and Sampling}
\label{sec:sampling}

\subsection{Viewer sample (n=400)}

The target population for Phase~1 was \emph{Tamil-film viewers
domiciled in Tamil Nadu and the Union Territory of Puducherry, aged
18 years and above, who had watched at least three Tamil-language
feature films in any mode in calendar year 2024}.  Although
population-frame data on Tamil-film viewers does not exist in a
listing form, two anchors were used to estimate the population
size: (i) the Tamil Nadu adult population of approximately 56
million (Census of India 2011 projected forward; \citealp{census2011,
ficciey2024}); and (ii) Ormax-reported penetration of Tamil-film
viewing of approximately 64\% of TN adults
\citep{ormax2023,ormax2024}, suggesting an effective population of
approximately 36 million.

\paragraph{Sample-size calculation.}  Using Cochran's
\citeyearpar{cochran1977} formula with a 95\% confidence interval,
margin of error of 5\% and conservative response variability
(\(p=0.5\)), the minimum sample size for an effectively infinite
population is 384 respondents.  We targeted 400 to allow for some
within-stratum imbalance and to round to a clean cell-design.

\paragraph{Sampling design.}  A \textbf{multi-stage stratified
sampling} approach was adopted:

\begin{enumerate}[itemsep=2pt]
\item \textbf{Stage 1 -- Geographic clusters.}  Eight geographic
      strata were defined: (i) Chennai (metropolitan); (ii)
      Coimbatore; (iii) Madurai; (iv) Tiruchirappalli; (v) Salem;
      (vi) Tirunelveli; (vii) other Tier-2 / Tier-3 (Erode, Vellore,
      Karaikudi, Tuticorin, Dindigul, Thanjavur, Kanyakumari,
      Karur); (viii) rural (selected blocks in each district).
      Allocation to strata reflects the urbanisation distribution
      of the Tamil-Nadu adult population, weighted modestly toward
      Chennai given its cinema-going salience.
\item \textbf{Stage 2 -- Within-cluster intercept and
      snowball.}  Within each cluster, respondents were recruited
      at five categories of intercept points -- multiplex foyers,
      single-screen theatres, college / university campuses, IT
      parks and shopping-mall food-courts -- supplemented by an
      online snowball recruitment via Tamil-cinema fan groups,
      WhatsApp networks and university alumni networks.
\item \textbf{Stage 3 -- Demographic-stratified quotas.}  Within
      each cluster, gender and age-band quotas (Table~\ref{tab:quotas})
      were enforced to approximate the Tamil-Nadu adult composition.
\end{enumerate}

\begin{table}[h]
\centering
\caption{Demographic quotas for the viewer sample (n=400)}
\label{tab:quotas}
\small
\begin{tabular}{lcc}
\toprule
\textbf{Stratum} & \textbf{Target \%} & \textbf{Achieved (n)}\\
\midrule
Age 18--24      & 27 & 112 \\
Age 25--34      & 33 & 136 \\
Age 35--44      & 21 & \phantom{0}80 \\
Age 45--54      & 13 & \phantom{0}48 \\
Age 55+         & \phantom{0}6 & \phantom{0}24 \\
\midrule
Female          & 45 & 176 \\
Male            & 54 & 220 \\
Other / no-answer  & \phantom{0}1 & \phantom{00}4 \\
\midrule
Chennai         & 26 & 104 \\
Coimbatore      & 14 & \phantom{0}56 \\
Madurai         & 11 & \phantom{0}44 \\
Tiruchirappalli & \phantom{0}9 & \phantom{0}36 \\
Salem           & \phantom{0}7 & \phantom{0}28 \\
Tirunelveli     & \phantom{0}6 & \phantom{0}24 \\
Other Tier-2/3  & 17 & \phantom{0}68 \\
Rural           & 10 & \phantom{0}40 \\
\bottomrule
\end{tabular}
\end{table}

\paragraph{Sample-size adequacy.}  An ex-post power analysis using
G*Power 3.1 \citep{faul2009}, assuming a medium effect size
$f^{2}=0.15$ for the multiple-regression analyses, $\alpha=0.05$ and
seven predictors, indicates statistical power above 0.99, well in
excess of the conventional 0.80 benchmark.

\subsection{Industry-specialist sample (n=50)}

The target population for Phase~2 was \emph{Tamil-cinema industry
professionals with at least three years' experience and at least one
credit / engagement on a film released in or after 2019}.  A
\textbf{purposive non-probability sample} of 50 specialists was
constructed to ensure role coverage across the value chain:

\begin{itemize}[itemsep=2pt]
  \item Producers -- 10
  \item Directors -- 8
  \item OTT acquisition / commissioning heads -- 7
  \item Distributors / sub-distributors -- 7
  \item Exhibitors (single-screen and multiplex) -- 6
  \item Marketing / PR professionals -- 5
  \item Music-label executives -- 4
  \item Trade analysts / journalists -- 3
\end{itemize}

Recruitment proceeded through a combination of direct outreach
(LinkedIn, professional email), industry-association referral
(through TFPC, TFAA and Producers Guild of India --- South Wing),
and snowball sampling from initial respondents.  All specialists
held informed-consent documentation and 32 of the 50 agreed to
audio-recorded interviews of 45--90 minutes; the remaining 18
completed only the structured instrument.  Twenty-eight interviews
were transcribed in full for thematic analysis.

\subsection{Generalisability}

The viewer sample is intended to support inference to the
\emph{Tamil Nadu / Puducherry adult Tamil-film-viewing population};
it is not intended to represent the diaspora.  The specialist
sample is purposive and supports analytical rather than statistical
generalisation \citep{yin2018}.  This is appropriate given the
small effective population of senior Tamil-cinema decision-makers
(estimated under 1{,}500 individuals).

\section{Research Instruments}
\label{sec:instruments}

\subsection{Viewer questionnaire (Appendix~A)}

The viewer questionnaire is divided into seven sections:

\begin{description}[style=nextline,leftmargin=2.4em]
\item[Section A: Demographics.]  Age, gender, location, education,
      household income.
\item[Section B: Theatre attendance.]  Pre-COVID and post-COVID
      monthly visits; preferred ticket band; primary screen format
      (single-screen vs multiplex).
\item[Section C: OTT subscription \& consumption.]  Number of
      subscriptions held; weekly hours; preferred platforms (Likert
      and multi-response).
\item[Section D: Genre, language and co-viewing.]  Pre vs post.
\item[Section E: Decision drivers.]  Star vs content priority;
      Likert items on word-of-mouth, music release, trailer, OTT
      recommendation, social-media virality, critic reviews.
\item[Section F: Willingness-to-pay.]  TVOD willingness-to-pay (INR
      bands); subscription churn behaviour; piracy self-report.
\item[Section G: Theatrical-window preference.]  Preferred wait
      between theatre and OTT release.
\end{description}

The instrument runs to 64 items, of which 41 are five-point Likert
items.  Five-point Likert was chosen following \citet{likert1932} and
\citet{tourangeau2000}; the mid-point ``Neither agree nor disagree''
was retained to avoid forced choice in culturally sensitive items.
The instrument was prepared in English and parallel-translated to
Tamil following the back-translation protocol of
\citet{brislin1970}, with translation conducted by two
bilingual research assistants and reconciled by the researcher.

\subsection{Specialist instrument and interview guide
(Appendix~B)}

The specialist instrument has two parts.  Part~I is a structured
questionnaire of 38 items covering: (a) respondent background;
(b) production-budget bracket distribution of recent projects;
(c) theatrical-window negotiation; (d) OTT-acquisition price
movement (indexed against 2018 baseline); (e) star-fee
restructuring; (f) marketing-channel mix; (g) revenue-mix
estimation; (h) employment shifts; (i) self-reported gratifications
and pain-points.  Part~II is a semi-structured interview guide of
22 open-ended prompts, organised in five clusters: pandemic
operating decisions; demand sense-making; capital-flow narratives;
labour-market dynamics; and three- to five-year outlook.

\section{Hypothesis Model}
\label{sec:hyp_model}

The eight hypotheses introduced in
Section~\ref{sec:hypotheses} of Chapter~\ref{ch:introduction} are
visualised in Figure~\ref{fig:hypothesis}.  The model treats the
COVID-19 disruption as the independent variable, OTT consumption and
theatrical demand as parallel mediating constructs, and industrial
reconfiguration as the dependent construct.

\begin{figure}[h]
\centering
\includegraphics[width=0.92\textwidth]{fig_34_hypothesis.png}
\caption{Hypothesis model with mediating constructs.}
\label{fig:hypothesis}
\end{figure}

\section{Pilot Study and Pre-test}
\label{sec:pilot}

A pilot of 40 viewer respondents and 6 industry specialists was
conducted in September 2024 in Chennai and Coimbatore.  The pilot
served four purposes: (i) to detect translation artefacts in the
Tamil version of the viewer questionnaire; (ii) to identify
respondent-burden hot-spots and trim items where possible; (iii) to
generate preliminary Cronbach's alpha estimates for Likert scales;
and (iv) to test the digital data-collection workflow (Google Forms
backed by an Airtable database).  Three questionnaire items were
revised post-pilot for clarity, two were dropped due to redundancy,
and one open-ended question was converted to a five-band
multiple-choice format.  The pilot data are not pooled with the
main analysis sample.

\section{Validity and Reliability}
\label{sec:reliability}

\paragraph{Content validity.}  Content validity was established
through a panel of five experts -- two academic supervisors, one
producer, one distributor and one trade analyst -- who rated each
item on a four-point relevance scale.  The Content Validity Index
(CVI) computed at the scale level (S-CVI/Ave) was 0.91, exceeding
the 0.80 threshold suggested by \citet{polit2007}.

\paragraph{Construct validity.}  Construct validity was assessed
through Exploratory Factor Analysis (EFA, principal-axis with
oblimin rotation) on the post-COVID Likert items.  Six factors
emerged with eigenvalues $>1$, jointly explaining 63.4\% of
variance.  Factor loadings on intended items exceeded 0.50 in all
but two cases; cross-loadings remained below 0.30.

\paragraph{Reliability.}  Cronbach's alpha was computed for each
multi-item scale (Table~\ref{tab:alpha}).  All scales exceeded the
0.70 threshold suggested by \citet{nunnally1994}, and most exceeded
0.80.

\begin{table}[h]
\centering
\caption{Reliability indices for primary multi-item scales}
\label{tab:alpha}
\small
\begin{tabular}{lccc}
\toprule
\textbf{Scale} & \textbf{Items} & \textbf{Cronbach's $\alpha$} & \textbf{KMO}\\
\midrule
OTT-gratification scale (post-COVID)            & 8 & 0.86 & 0.83\\
Theatre-experience valuation scale              & 6 & 0.81 & 0.79\\
Decision-driver scale (star vs content)         & 7 & 0.78 & 0.77\\
Genre-preference scale (post-COVID)             & 10& 0.83 & 0.81\\
Willingness-to-pay scale (TVOD)                 & 5 & 0.74 & 0.72\\
Specialist outlook scale                        & 9 & 0.84 & 0.78\\
\bottomrule
\end{tabular}
\end{table}

\paragraph{Triangulation.}  Inter-method triangulation between
viewer survey, specialist responses and trade-data sources is the
fourth quality-control layer.  Where a finding is supported by all
three sources it is treated as robust; where it is supported by
only one or two, the discussion in Chapter~\ref{ch:conclusion}
flags this.

\section{Data-Collection Procedure}
\label{sec:datacoll}

Data collection ran in three field phases:

\begin{itemize}[itemsep=2pt]
\item \textbf{October 2024 -- February 2025.}  Viewer survey
      across the eight geographic strata.  Hybrid mode of
      administration: 224 respondents completed in person on a
      tablet (Google Forms), 176 completed online via a circulated
      link.  Median completion time was 19 minutes.  Response
      validation rules in the form (range checks, mandatory items,
      consistency rules) ensured data integrity at point of entry.
\item \textbf{March 2025 -- May 2025.}  Specialist questionnaire
      and interviews.  All structured-instrument data were captured
      digitally; interviews were audio-recorded with consent and
      transcribed verbatim by a professional transcription
      service, with researcher-led verification of all
      transcripts.
\item \textbf{June 2025 -- July 2025.}  Secondary data extraction
      from FICCI--EY 2018--2024 reports, KPMG 2019--2024, BCG
      2018--2023, Ormax 2022--2024, Statista 2019--2024, MPA
      2023--2024 and a curated archive of trade-press articles.
\end{itemize}

\section{Data-Analysis Techniques}
\label{sec:analysis}

\subsection{Quantitative analysis}

All quantitative analyses were performed in IBM SPSS Statistics 28,
with selected models replicated in R 4.4.0.  The analytical
techniques are mapped to research questions in
Table~\ref{tab:rq_method}.

\begin{table}[h]
\centering
\caption{Mapping of research questions to analytical techniques}
\label{tab:rq_method}
\small
\begin{tabular}{p{3.5cm} p{4.0cm} p{6.0cm}}
\toprule
\textbf{Research Question} & \textbf{Hypothesis(-es)} & \textbf{Technique}\\
\midrule
RQ1 (viewership change)  & H1, H2, H3, H4 &
Descriptive statistics; paired-sample $t$-test; Wilcoxon signed-rank;
McNemar test; chi-square test of homogeneity. \\
RQ2 (demographic moderators) & H5 &
One-way and factorial ANOVA; multiple regression; moderated regression. \\
RQ3 (industry reconfiguration) & H6, H7, H8 &
Descriptive statistics on specialist data; one-sample $t$-tests vs
2018 baseline; thematic analysis of interview data. \\
RQ4 (integrated framework) & --- &
Mixed-method synthesis; framework induction. \\
\bottomrule
\end{tabular}
\end{table}

\subsection{Qualitative analysis}

Interview transcripts were analysed in QSR NVivo 14 using a
\textbf{constant-comparative} approach
\citep{strauss1998,saldana2021}.  A first-cycle of descriptive and
in-vivo coding produced 247 codes; a second-cycle of pattern coding
collapsed these to 38 categories; a third cycle of theoretical
coding produced 9 themes that map onto the Hybrid-Equilibrium
framework presented in Chapter~\ref{ch:conclusion}.  Inter-coder
agreement was assessed on a 20\% sub-sample by an independent
research assistant; Cohen's kappa was 0.81, indicating substantial
agreement \citep{landis1977}.

\subsection{Mixed-method integration}

Mixed-method integration was achieved through (a) a
\emph{joint-display} matrix \citep{guetterman2015} that arrays
quantitative findings against qualitative themes; and (b) a
\emph{narrative-weave} discussion in Chapter~\ref{ch:conclusion}
that interleaves quantitative effect sizes with illustrative
specialist quotations.

\section{Ethical Considerations}
\label{sec:ethics}

The research was reviewed and approved by the Institutional Ethics
Committee of \institution\ (Approval number IEC/2024/Media-XX, dated
[date]).  The following ethical safeguards were implemented:

\begin{itemize}[itemsep=2pt]
\item \textbf{Informed consent.}  All viewer respondents read /
      were read a written consent form before data collection.
      All specialists signed a consent form prior to interview.
\item \textbf{Anonymity and confidentiality.}  Viewer responses
      are reported only in aggregate.  Specialist quotations are
      attributed only by role (``a senior OTT acquisition head'',
      ``a multiplex chain executive'') unless explicit consent for
      named attribution was given.
\item \textbf{Data security.}  All raw data are stored in
      access-controlled cloud storage; interview audio files are
      encrypted at rest and will be deleted at the end of the
      mandated five-year retention period.
\item \textbf{No sensitive personal data.}  No data on caste,
      religion, sexual orientation or political affiliation were
      sought.
\item \textbf{Right to withdraw.}  All participants were informed
      of their right to withdraw at any point without consequence;
      two viewer respondents and one specialist exercised this
      right and their data were deleted.
\item \textbf{Compensation.}  Viewer respondents received a
      nominal mobile-recharge voucher (INR 50) on completion;
      specialists received no monetary compensation.
\end{itemize}

\section{Chapter Summary}
\label{sec:meth_summary}

This chapter has set out the pragmatist-mixed-methods design of the
study, the multi-stage stratified sampling of 400 viewers and the
purposive sampling of 50 industry specialists, the construction and
psychometric properties of the two instruments, the eight-hypothesis
testing model, the data-collection workflow, the SPSS- and
NVivo-based analytical techniques and the ethical safeguards.  The
chapter has also justified each design choice against alternative
options.  The next chapter situates the empirical analysis by
providing a historical and contemporary overview of the Tamil film
industry against which the primary data are interpreted.
